\chapter{遗忘的计算动力学：解绑定、降权与结构回收}
\label{chap:geometry_of_oblivion}
遗忘不是单一的“信息消失”，也不能由退相干、能量耗散或曲率平滑直接证明。一个主体可能暂时检索不到名称，却仍保有语义线索；可能记得事件事实，却不再具有当时的情绪反应；也可能因索引迁移、权限改变、模型压缩或存储删除而表现出不同类型的缺失。本章把遗忘拆成身份—内容解绑定、访问概率下降、细节压缩、价值再估计、结构剪枝和授权删除等可区分过程，研究它们在三相记忆与外部记忆之间如何发生、何时有助于控制复杂度、何时造成不可恢复损失。物理语言只保留为辅助意象，所有机制以对象、版本、时间尺度、证据和任务后果验收。

\section{记忆对象与遗忘操作的类型化分解}
\label{sec:tensor_definition_of_memory}
原章把完整记忆定义为“形算子、质矢量与相位”的张量积，再把遗忘写成张量积依次退化为直和、零向量和平坦度量。这个表示没有说明对象身份、来源、时间、访问权限和读出任务，也无法区分“内容仍在但暂时找不到”与“内容已经被删除”。我们从可追溯记录开始。令单个情景记忆为
\[
m_j=(\iota_j,x_j,c_j,t_j,v_j,s_j,p_j,a_j,r_j),
\]
其中 $\iota_j$ 是稳定身份键，$x_j$ 是内容表示，$c_j$ 是人物、地点与任务上下文，$t_j$ 是带置信度的时间信息，$v_j$ 是价值或情绪估计，$s_j$ 是来源谱系，$p_j$ 是访问权限，$a_j$ 是当前可访问度，$r_j$ 是保留策略。各分量类型不同：内容向量可在表示空间中连续更新，权限是离散规则，来源是图结构，时间是索引及不确定区间，不能用一个“记忆强度”替代。

记忆是否能在任务 $q$ 下被读出，还取决于查询、索引和当前状态。设检索器返回候选分布
\[
P_k(j\mid q,h_k)=\frac{\exp S(q,h_k,m_j)}{\sum_{\ell\in\mathcal A_k}\exp S(q,h_k,m_\ell)},
\]
其中 $S$ 是经过尺度校准的匹配分数，$\mathcal A_k$ 是权限与有效期允许访问的集合。即使 $m_j$ 未变化，查询线索、竞争记录、索引版本或权限集合改变，也会使返回概率下降；这属于访问性遗忘，不等于存储消失。反过来，高匹配分数只支持候选检索，不证明内容真实。读出后仍须检查身份、来源冲突与环境证据。

我们把遗忘操作分为六类。第一，解绑定：身份键与内容、人物或来源链接变弱或冲突，典型表现是知道某个词的语义和首字母却无法恢复完整名称。第二，访问降权：记录仍在，但其索引优先级或任务相关性下降。第三，细节压缩：保存事件摘要而丢弃感官或时序细节。第四，价值再估计：事实内容仍可恢复，而情绪或风险读出随新经验改变。第五，结构剪枝：跨事件形成的慢结构不再获得证据支持，被候选替代或降级。第六，授权删除：依照隐私、法律或治理规则移除记录及派生索引。这些过程可以同时发生，却不构成必然的先后相序。

在三相记忆中，$h(t)$ 的激活消退通常只是当前工作态释放；$E$ 中事件的降权、压缩和解绑定改变未来检索；$\Theta$ 的结构回收改变预测与策略；外部记忆Me还可能保留内部系统已不再访问的原始证据。因而“我想不起来”可能源于当前容量不足、线索不匹配、权限禁止、索引错误、版本冲突或真实删除。诊断必须返回失败类型，而不能统一报告为空值。AgentOS可以用CCM降低低收益记录在$h(t)$中的竞争，用Boundary治理删除和外部访问，用Offloading把可追溯原件移入Me；这些只是工程实例，并非所有智能系统的唯一机制。

原章提出“有形无质”和“有质无形”，我们保留其中的组合问题，但改写为类型化读出。前者可能是语义、人物关系或首字母等结构线索尚存，而语音形式未被检索；后者可能是熟悉度或价值读出升高，却缺乏可靠的事件来源绑定。舌尖现象不能证明纤维处于基态，既视感也不能证明游离孤立子存在。可检验的做法是操纵提示类型、候选竞争和延迟，分别测量部分信息、命中率、错误来源与主观置信，并比较结构线索和内容恢复是否可独立变化。

离散实现必须保留版本。对任一操作 $o_k$，系统先生成候选 $m'_j=o_k(m_j)$，记录原因、授权者、输入证据和受影响索引，再在影子副本上复演关键任务。只有满足保留、隐私、安全和可恢复性条件才提交；高风险删除应先建立加密或隔离备份，并声明恢复期限。对人类记忆，本章的计算分类只能组织问题，不能据此诊断失忆、创伤或既视感。验收包括同义线索、错误人名、密码更新、情绪淡化、索引迁移、权限撤销与真实删除，报告检索概率、来源准确、压缩误差、旧任务退化和恢复成本。这样，遗忘首先成为一组可区分操作，而不是由张量符号或物理名称预先决定的三段命运。
